% SC1008 — Chapter 8: Modern C++ — RAII, Smart Pointers, Move (0->100 ground-up)
\chapter{Modern C++: RAII, Smart Pointers, and Move Semantics}
\label{ch:modern-cpp}

\begin{quote}
\emph{``Modern C++ is not C with extras --- it is a different language that happens to compile C.''} RAII makes resource safety automatic; smart pointers make ownership explicit; move semantics makes copies cheap.
\end{quote}

\noindent\fbox{\parbox{\dimexpr\linewidth-2\fboxsep-2\fboxrule}{%
\textbf{From zero: Ch.~4's discipline was manual; now we automate it.} Every resource gets an owner object whose destructor frees it. We build RAII, then smart pointers, then moves --- each fixing a failure mode of the last.}}

\section*{Learning Outcomes}
After this chapter you will be able to:
\begin{enumerate}[label=(L\arabic*)]
    \item explain RAII and why it prevents leaks even on exceptions;
    \item use \texttt{unique\_ptr}, \texttt{shared\_ptr}, \texttt{weak\_ptr} correctly;
    \item distinguish lvalues from rvalues and apply move semantics;
    \item state the Rule of Zero/Three/Five;
    \item write exception-safe code with RAII types.
\end{enumerate}

% ============================================================
\section{RAII: Resource Acquisition Is Initialisation}
\label{sec:raii}
\emph{Idea:} tie every resource (memory, file, lock) to an object's lifetime. Constructor acquires; destructor releases. No manual \texttt{free}/\texttt{close} --- the scope does it, even if an exception is thrown.

\begin{figure}[h]\centering
\begin{tikzpicture}[scale=0.68,every node/.style={font=\scriptsize}]
\draw[thick,fill=green!12] (0,0) rectangle (3.0,1.1);
\node[font=\tiny\bfseries] at (1.5,0.85) {Constructor};
\node[font=\tiny] at (1.5,0.50) {acquire (new/open/lock)};
\node[font=\tiny] at (1.5,0.20) {\texttt{File f("a.txt")}};
\draw[thick,fill=blue!10] (3.6,0) rectangle (6.6,1.1);
\node[font=\tiny\bfseries] at (5.1,0.85) {Use};
\node[font=\tiny] at (5.1,0.50) {read/write};
\node[font=\tiny] at (5.1,0.20) {exception? still OK};
\draw[thick,fill=orange!15] (7.2,0) rectangle (10.2,1.1);
\node[font=\tiny\bfseries] at (8.7,0.85) {Destructor};
\node[font=\tiny] at (8.7,0.50) {release (delete/close/unlock)};
\node[font=\tiny] at (8.7,0.20) {automatic at \texttt{\}}};
\draw[-{Latex},thick,blue!60!black] (3.0,0.55)--(3.6,0.55);
\draw[-{Latex},thick,blue!60!black] (6.6,0.55)--(7.2,0.55);
\draw[-{Latex},thick,red!60!black,dashed] (5.1,-0.05)--(5.1,-0.45)--(8.7,-0.45)--(8.7,0.0);
\node[font=\tiny,red!60!black] at (6.9,-0.70) {exception $\to$ stack unwind $\to$ dtor still runs};
\end{tikzpicture}
\caption{RAII timeline: acquisition in constructor, automatic release in destructor. Stack unwinding guarantees the destructor runs even on exception.}
\label{fig:raii-timeline}
\end{figure}

\begin{lstlisting}[language=C++]
#include <fstream>
void process(){
    std::ofstream out("log.txt"); // acquires file handle
    out << "hello\n";
    if(some_error()) throw std::runtime_error("oops");
    // no close needed -- destructor closes even on throw
} // out.~ofstream() runs here
\end{lstlisting}
\texttt{std::string}, \texttt{std::vector}, \texttt{std::lock\_guard}, \texttt{std::fstream} are all RAII types. Manual \texttt{new/delete} is the non-RAII alternative --- avoid it.

\section{Smart Pointers: Explicit Ownership}
\label{sec:smart-pointers}

\begin{figure}[h]\centering
\begin{tikzpicture}[scale=0.68,every node/.style={font=\scriptsize}]
% unique_ptr
\draw[thick,fill=blue!10] (0,1.2) rectangle (2.4,1.95);
\node[font=\tiny\bfseries] at (1.2,1.62) {\texttt{unique\_ptr<T> p}};
\draw[thick,fill=green!14] (3.2,1.2) rectangle (5.0,1.95);
\node at (4.1,1.57) {\texttt{T object}};
\draw[-{Latex},thick,red!70!black] (2.4,1.57)--(3.2,1.57) node[midway,above,font=\tiny] {owns};
\node[font=\tiny,fill=yellow!12,rounded corners=2pt,inner sep=2pt] at (1.2,0.85) {one owner, move-only};
% shared_ptr
\draw[thick,fill=blue!10] (0,-0.15) rectangle (2.0,0.55);
\node[font=\tiny\bfseries] at (1.0,0.22) {\texttt{shared\_ptr q}};
\draw[thick,fill=blue!10] (2.6,-0.15) rectangle (4.6,0.55);
\node[font=\tiny\bfseries] at (3.6,0.22) {\texttt{shared\_ptr r}};
\draw[thick,fill=green!14] (5.4,-0.15) rectangle (7.2,0.55);
\node at (6.3,0.20) {\texttt{T}};
\draw[thick,fill=orange!15] (5.4,-0.95) rectangle (7.2,-0.35);
\node[font=\tiny] at (6.3,-0.65) {control block: ref=2};
\draw[-{Latex},thick,red!70!black] (2.0,0.22)--(5.4,0.22);
\draw[-{Latex},thick,red!70!black] (4.6,0.22)--(5.4,0.22);
\draw[thick,dashed,gray] (5.4,0.55)--(5.4,-0.35);
\node[font=\tiny,fill=yellow!12,rounded corners=2pt,inner sep=2pt] at (1.0,-0.95) {shared: ref-counted};
\node[font=\tiny,fill=violet!10,rounded corners=2pt,inner sep=2pt] at (3.6,-0.95) {\texttt{weak\_ptr} breaks cycles};
\end{tikzpicture}
\caption{Smart pointers: \texttt{unique\_ptr} has one owner; \texttt{shared\_ptr} shares via a control block with reference count; \texttt{weak\_ptr} observes without owning.}
\label{fig:smart-ptrs}
\end{figure}

\begin{lstlisting}[language=C++]
#include <memory>
auto p = std::make_unique<int>(42);        // unique ownership
// auto q = p;                             // error: not copyable
auto q = std::move(p);                     // transfer ownership; p==nullptr

auto a = std::make_shared<std::string>("hello"); // shared, ref count 1
auto b = a;                                // ref count 2, same object
std::cout<<a.use_count()<<"\n";            // 2
std::weak_ptr<std::string> w = a;          // weak: doesn't increment count
// w.lock() -> shared_ptr if object still alive
\end{lstlisting}
Prefer \texttt{make\_unique}/\texttt{make\_shared} over raw \texttt{new}. Use \texttt{unique\_ptr} by default; \texttt{shared\_ptr} only when shared ownership is genuinely needed. \texttt{weak\_ptr} breaks cycles (\texttt{A} owns \texttt{B}, \texttt{B} observes \texttt{A} weakly).

\section{Move Semantics: Stealing Instead of Copying}
\label{sec:move}
Copying a vector duplicates its heap buffer ($O(n)$). \emph{Moving} steals the buffer pointer and nulls the source ($O(1)$). Moves fire on \emph{rvalues} (temporaries, \texttt{std::move}'d values).

\begin{figure}[h]\centering
\begin{tikzpicture}[scale=0.68,every node/.style={font=\scriptsize}]
% copy
\node[font=\tiny\bfseries] at (1.6,1.35) {Copy: duplicate buffer};
\draw[thick,fill=blue!10] (0,0.45) rectangle (1.6,1.05);
\node[font=\tiny] at (0.8,0.75) {\texttt{a}};
\draw[thick,fill=green!12] (0.3,0.0) rectangle (1.3,0.35);
\node[font=\tiny] at (0.8,0.17) {buf};
\draw[thick,fill=blue!10] (2.2,0.45) rectangle (3.8,1.05);
\node[font=\tiny] at (3.0,0.75) {\texttt{b}};
\draw[thick,fill=green!12] (2.5,0.0) rectangle (3.5,0.35);
\node[font=\tiny] at (3.0,0.17) {buf copy};
\draw[-{Latex},thick,blue!60!black] (1.6,0.75)--(2.2,0.75) node[midway,above,font=\tiny] {$O(n)$};
% move
\node[font=\tiny\bfseries] at (6.8,1.35) {Move: steal pointer};
\draw[thick,fill=blue!10] (5.2,0.45) rectangle (6.8,1.05);
\node[font=\tiny] at (6.0,0.75) {\texttt{a (moved-from)}};
\draw[thick,fill=gray!12] (5.5,0.0) rectangle (6.5,0.35);
\node[font=\tiny,gray] at (6.0,0.17) {null};
\draw[thick,fill=blue!10] (7.4,0.45) rectangle (9.0,1.05);
\node[font=\tiny] at (8.2,0.75) {\texttt{b}};
\draw[thick,fill=green!14] (7.7,0.0) rectangle (8.7,0.35);
\node[font=\tiny] at (8.2,0.17) {buf};
\draw[-{Latex},thick,red!70!black] (6.8,0.75)--(7.4,0.75) node[midway,above,font=\tiny] {$O(1)$};
\draw[-{Latex},thick,red!60!black,dashed] (6.0,0.45)--(8.2,0.45);
\end{tikzpicture}
\caption{Copy duplicates the heap buffer; move transfers the pointer, leaving the source empty. Move is $O(1)$ vs.\ $O(n)$.}
\label{fig:move-vs-copy}
\end{figure}

\begin{lstlisting}[language=C++]
std::vector<int> a = {1,2,3,4,5};
std::vector<int> b = a;                // copy: O(n)
std::vector<int> c = std::move(a);     // move: O(1), a is now empty
// a is valid but unspecified -- only safe to assign/destroy
std::vector<int> make_vec(){ return {1,2,3}; } // NRVO/move elision
auto d = make_vec();                   // no copy (elided)
\end{lstlisting}
The compiler auto-moves temporaries; write \texttt{std::move} only for named lvalues you no longer need. After move, treat the source as empty.

\section{Rule of Zero / Three / Five}
\label{sec:rule-of-five}
\begin{itemize}[leftmargin=*]
    \item \textbf{Rule of Zero:} if your class uses RAII members (\texttt{string}, \texttt{vector}, smart pointers), write \emph{no} custom destructor/copy/move --- the compiler does it correctly.
    \item \textbf{Rule of Five:} if you \emph{must} write one of destructor, copy-ctor, copy-assign, move-ctor, move-assign, you likely need all five.
    \item Prefer Zero; fall to Five only for low-level resource wrappers.
\end{itemize}
\begin{lstlisting}[language=C++]
// Rule of Zero: no custom special members needed
class Student{
    std::string name_; std::vector<int> scores_;
public:
    Student(std::string n): name_(std::move(n)) {}
    // compiler-generated dtor/copy/move are correct
};
\end{lstlisting}


\section{Exception Safety and RAII in Practice}
RAII makes exception safety almost free. Consider manual vs.\ RAII resource handling when an exception is thrown:

\begin{lstlisting}[language=C++]
void manual(){
    int *p = new int[100];
    FILE *f = fopen("data.txt","r");
    risky_operation(); // if throws, leak p and f!
    delete[] p; fclose(f);
}
void raii(){
    auto p = std::make_unique<int[]>(100);
    std::ifstream f("data.txt"); // RAII file handle
    risky_operation(); // if throws, destructors still run -- no leak
} // p and f destroyed automatically
\end{lstlisting}
Exception safety levels: \emph{basic} (no leak), \emph{strong} (rollback on failure), \emph{nothrow} (never throws). RAII types give basic safety by default; transactions give strong safety.

\section{Perfect Forwarding and Modern Idioms}
C++11 introduced \texttt{auto}, range-for, lambdas, and forwarding references. C++17/20 added structured bindings and concepts.

\begin{lstlisting}[language=C++]
auto add = [](auto a, auto b){ return a+b; }; // generic lambda (C++14)
std::cout<<add(3,4)<<" "<<add(3.1,2.9)<<"\n";  // 7 6.0

std::map<std::string,int> m={{"Asha",95},{"Bo",82}};
for(auto &[name,score]: m) // structured binding (C++17)
    std::cout<<name<<":"<<score<<"\n";

// Perfect forwarding: preserve lvalue/rvalue nature
template<typename T>
void wrapper(T&& arg){ // forwarding reference
    callee(std::forward<T>(arg)); // forwards as l- or r-value
}
\end{lstlisting}
These idioms reduce boilerplate and let the compiler generate optimal code. Prefer \texttt{auto} when the type is obvious; prefer range-for over index loops; prefer lambdas over hand-written functors.

\section{Chapter Notes}
Meyers \emph{Effective Modern C++} (2014) is the definitive move/RAII guide. C++ Core Guidelines: ``R.1: Manage resources automatically using RAII.''

\section{Exercises}
\begin{enumerate}[label=E8.\arabic*, leftmargin=*, itemsep=0.6em]
\item \textbf{RAII proof.} Write a function that opens two files and throws midway. Show RAII handles close it correctly while manual \texttt{fopen/fclose} leaks the second file.
\item \textbf{Ownership.} When would you choose \texttt{unique\_ptr} vs.\ \texttt{shared\_ptr} vs.\ raw \texttt{T*} (non-owning)? Give an example of each.
\item \textbf{Move trace.} \texttt{vector<int> a=\{1,2\}; vector<int> b=std::move(a); cout<<a.size()<<" "<<b.size();} What prints? Is \texttt{a[0]} safe to read after the move?
\item \textbf{Rule of Zero.} A class holds \texttt{string name\_} and \texttt{vector<double> grades\_}. Do you need a custom destructor? What if it held \texttt{int* buf} from \texttt{new int[10]}?
\item \textbf{Cycle.} \texttt{struct A\{shared\_ptr<B> b;\}; struct B\{shared\_ptr<A> a;\};} Why does this leak? Fix it with \texttt{weak\_ptr} and explain the ref counts.
\end{enumerate}
% End of Chapter 8
